\section{Introduction}
\label{sec:intro}

The summed neutrino mass $\mnu$ is one of the main targets of precision cosmology. Its constraints come from combining early-universe information in the cosmic microwave background (CMB) with late-time distance measurements such as baryon acoustic oscillations (BAO). As upper limits approach the minimum allowed by oscillation experiments, a practical question becomes pressing: when a limit moves, how much of the movement reflects the sky, and how much reflects choices in how the data were processed and compressed into a likelihood?

A sharp example is reported in \cite{arxiv2601_16277}. The authors compare two releases of the South Pole Telescope SPT-3G temperature and polarization (TT/TE/EE) likelihood: the 2018 release and the later D1 release built from the 2019--2020 seasons. Holding the rest of their CMB baseline fixed (a subset of Planck 2018 products plus Planck PR4 lensing, together with the matching SPT lensing product), they find that switching from the 2018 to the D1 combination tightens the $\mnu$ upper limit and moves the posterior towards the physical boundary $\mnu=0$. They call the shift ``curious'', because both releases observe essentially the same field, and they point out that the decisive test would be to reprocess the 2018 raw data through the D1 pipeline. We do not claim that test as our idea, and we do not perform it.

\paragraph{The question we ask.}
Before interpreting a parameter shift, one can ask whether the two released likelihoods agree with each other at all. We frame this with the Six Birds framework \cite{sixbirds_foundations,sixbirds_darkenergy}. Each likelihood release is a \emph{lens}: a map from the sky (together with instrument and foreground nuisances) to a likelihood object. Here ``lens'' never means gravitational lensing. A model class with its nuisance parameterization and priors is a \emph{completion}. The posterior we report is the \emph{package} obtained by composing the two. If two lenses describe the same sky consistently, parameters chosen by one should fit the other reasonably well. We measure how far this fails, in each direction separately, and where in the data the failure sits.

\paragraph{What this paper contributes.}
\begin{enumerate}[leftmargin=*, itemsep=2pt]
  \item \textbf{Directional cross-lens audits.} We fit one SPT release (with a common completion), transfer the shared cosmological parameters to the other release, and record the loss in fit relative to that release's own best fit. We do this in both directions, for two baselines (SPT with DESI DR2 BAO, and SPT with a Planck subset, Planck PR4 lensing and DESI), and under three controls: a multipole-support cut, removal of the SPT optical-depth prior, and higher Boltzmann precision (\cref{sec:results-audits}).
  \item \textbf{Exact localization with correlated data.} We split the SPT quadratic part of each mismatch over spectra and multipole bands with a signed allocation that uses the full covariance matrix. The pieces sum exactly to that quadratic part, unlike block-by-block chi-squares, which fail to add up when blocks are correlated; native prior and normalization terms are accounted for separately (\cref{sec:ledger}). The resulting maps show that the location of the mismatch depends on both the baseline and the direction of transfer.
  \item \textbf{A clear statement of what the $\mnu$ chains support.} Our sampled posteriors do not pass convergence and quantile-precision checks fixed in advance, and the likelihood evaluations turned out to be sensitive to solver version and numerical settings. We therefore report no magnitude for the $\mnu$ tightening, and we explain what a qualified claim would require (\cref{sec:results-mnu}).
  \item \textbf{A verified boundary mechanism.} For two Gaussian constraints and the gate $\mnu\ge0$, we state exactly when the posterior mode sits at zero, show that two constraints that both prefer positive masses can never put it there, and show that tightening one constraint can make a boundary mode either more or less likely depending on the other constraint's sign. All statements are checked in Lean~4 with Mathlib \cite{lean4,mathlib} (\cref{sec:boundary-mode} and Appendix~\ref{app:formal}).
\end{enumerate}

\paragraph{What this paper does not claim.}
We do not identify the instrumental or analysis cause of the disagreement, and we do not claim a correct value or bound for $\mnu$. Our audit values come from numerical optimizations that are not certified to be global optima, and they are not calibrated significance tests. The SBT language provides the vocabulary of the audits. It is not a theorem that supplies the missing physical steps.

\paragraph{Outline.}
\Cref{sec:related-work} places the work in context. \Cref{sec:sbt-framework} introduces lenses, packages and the directional audit. \Cref{sec:baselines} defines the data combinations and \cref{sec:methods} the audit and sampling procedures. \Cref{sec:results-audits} presents the audits and \cref{sec:results-mnu} the status of the $\mnu$ posteriors. \Cref{sec:boundary-mode} analyses boundary modes, and \cref{sec:discussion,sec:conclusion} discuss limitations and conclude. The appendices list the formal results, describe reproducibility, and record the version history.
